\section{Introduction}
\label{Intro}

Multiple Control Toffoli (MCT) gates are essential parts of quantum algorithms. Coming from
reversible computing, they are also used in quantum computing for large, non-linear
entangling operations.

MCT are part of a set of operations known as non-Clifford~\cite{bib-barenco-elementary}
operations. As MCT gates are not elementary gates, they are constructed using a basis gate set,
often using the Clifford+T set, a uniersal set of quantum gates~\cite{bib-bravyi2005universal}.
Amy et al proved that there is a minimum number of T-gates,
which implement non-Clifford operations, to implement the MCT~\cite{amy-meet-in-middle}.
T-gates are especially important because in fault-tolerant quantum computing, they have a
higher implementation cost than other gates in the Cliffford+T set~\cite{bib-bravyi2005universal}.
The number of T-gates, as well as their inverse the T$^{\dagger}$ gate, in a quantum
circuit is known as the \emph{T-count}. The number of such gates in a dependency chain is known
as its \emph{T-depth}. T-gates dominate the implementation of a quantum
circuit so much, optimizing the T-count, especially for large components such as the MCT, is
an imperative for quantum circuit optimization.

Several techniques have been developed to reduce the T-count. A technique to reduce the number
of non-Clifford rotations in a quantum circuit was introduced by Barenco et al~\cite{bib-barenco-elementary},
which allows a \emph{relative-phase}, or an input dependent phase, on the outputs. Other techniques,
such as \emph{clean ancilla} (initialized to 0)~\cite{bib-mike-and-ike} and \emph{dirty ancilla}
(no restrictions on value)~\cite{bib-barenco-elementary} decompositions mitigate the asymptotic T-count
of the MCT.

In general, quantum circuits do not allow for a relative phase, as it can affect other inputs downstream.
Other work such as from Maslov~\cite{bib-maslov-rtof} and Amy~\cite{bib-amy-phase-state}, attempt
to combine relative-phase implementations and the clean ancilla decompositions to reduce the cost.
This is because the Toffoli gates that comprise the decompositions can themselves be implemented
up to relative phase, because there is a need to \emph{uncompute}, or restore the ancilla's value,
them later. This uncomputation logic can then erase the relative phase. Other works such
as Zindorf et al~\cite{bib-zindorf-mct} and Abdessaied et al~\cite{bib-abdessaied-mct} instead use
the dirty ancilla decomposition. Zindorf et al in particular achieves the best known T-count for
the MCT using only one dirty ancilla. 

In this work, we employ sum-of-paths techniques used in~\cite{bib-amy-phase-state} and~\cite{bib-amy-cnot}
to analyze existing techniques and develop gadgets to further reduce the T-count in these MCT
gates. We find that by employing relative phase techniques on dirty ancilla-like constructions such
as those used in~\cite{bib-zindorf-mct}, we can achieve T-counts equal to or almost equal to the
clean ancilla construction with only a single dirty ancilla. This is the best known T-count for MCT
to the authors' knowledge.

This paper is structured as follows. In~\secref{Pre}, we first introduce preliminary
knowledge to understand our proposals, including MCT decomposition and
relative phase. We then present our proposal in~\secref{Pro}, including the relative
phase dirty ancilla synthesis. Finally, \secref{Conc} concludes with our conclusions and future work.
